\chapter{Pfarrfrauen}
\quellseite{9}

Aus den Neu-England-Staaten kam vor einiger Zeit die Kunde, daß sich mehrere Pastoren von ihren Frauen haben scheiden lassen. In jedem einzelnen Fall wurde als Grund angegeben, die Frau habe den Mann verlassen. Dieses wiederum wurde von Seiten der Verklagten damit zu rechtfertigen gesucht, daß es ihnen bei dem geringen Einkommen ihres Mannes unmöglich gewesen wäre, ihren gesellschaftlichen Verpflichtungen nachzukommen, und daß es überhaupt langweilig sei, eine Pfarrfrau zu sein. Diese Frauen waren offenbar Modedamen; sie hatten keinen Begriff von dem Predigtamt und ihrer Stellung; sie gehörten nicht in das Pfarrhaus, darum haben sie es auch treulos verlassen. Das Wirken der Pfarrfrauen wird vielfach sehr gering geschätzt. Gar manche von ihnen sind Heldinnen, Heldinnen in der großen Kunst der Selbstverleugnung, der Entsagung. Bei einem geringen Einkommen und einer häufig nicht kleinen Familie den Haushalt im Geleise zu halten, all die verschiedenen Anforderungen, die an sie gestellt werden, zu befriedigen und dabei nicht den Muth zu verlieren, sondern sogar den Frohsinn zu bewahren, wahrlich, das erfordert eine Charakterstärke, eine Glaubenszuversicht, die die Welt nicht zu würdigen weiß. Diese rechten Pfarrfrauen, die der Pfarrer Stütze und Aufmunterung sind, sind Lichter, die im Verborgenen leuchten und sich verzehren, indem sie leuchten. Gott segne alle die Pfarrfrauen, die ihren Beruf recht verstehen! Am jüngsten Tage wird man einmal sehen, was sie gethan haben.

(Friedensbote.)
